% Source: Solutions/Yash-Indrajit-Himanshu-PS6.pdf, p. 2 (typed, joint submission; Q1-3 credited to Yashpal Prasad).
\submittedby{Yashpal Prasad}

\begin{proof}
Let $p\colon X\to Y$ and $q\colon Y\to Z$ be quotient maps. By definition, both $p$ and $q$ are
continuous and surjective maps. Consider the composite map $qp\colon X\to Z$ defined by
$(qp)(x)=q(p(x))$. The composition of continuous functions is continuous, and the composition of
surjective functions is surjective; hence $qp$ is continuous and surjective.

To prove that $qp$ is a quotient map, we must show that a subset $U\subseteq Z$ is open in $Z$
if and only if its inverse image $(qp)^{-1}(U)$ is open in $X$. By set-theoretic properties of
inverse images under compositions, $(qp)^{-1}(U)=p^{-1}(q^{-1}(U))$.

Suppose $U\subseteq Z$ is such that $(qp)^{-1}(U)$ is open in $X$. Since $p$ is a quotient map,
a subset $V\subseteq Y$ is open in $Y$ if and only if $p^{-1}(V)$ is open in $X$. Setting
$V=q^{-1}(U)$, we have $p^{-1}(V)=(qp)^{-1}(U)$ open in $X$, which implies $q^{-1}(U)$ is open in
$Y$. Now, since $q$ is a quotient map and $q^{-1}(U)$ is open in $Y$, it follows that $U$ is
open in $Z$. Conversely, if $U$ is open in $Z$, then $q^{-1}(U)$ is open in $Y$ by continuity
of $q$, and $p^{-1}(q^{-1}(U))$ is open in $X$ by continuity of $p$.

\textbf{Conclusion:} $U$ is open in $Z$ if and only if $(qp)^{-1}(U)$ is open in $X$. Thus,
$qp$ is a quotient map.
\end{proof}
